\begin{abstract}
Large language model (LLM) agents that integrate external tools are vulnerable to indirect prompt injection (IPI), where malicious instructions embedded in tool-returned content hijack the agent into executing attacker-specified actions. A growing body of work proposes defenses---input encoding, structured queries, runtime task alignment, attention-based detection, and architecture-level isolation---yet these defenses are typically evaluated only against static injection templates. Recent work by Zhan et al.\ showed that adaptive attacks, which tailor the injection to the defense's detection signal, break three representative defenses. We extend this line of inquiry to its logical conclusion: we systematize 12 IPI defenses across five defense classes, build a unified adaptive attack framework that materializes a matched adversary for each class, and conduct a large-scale evaluation on the AgentDojo benchmark across 3{,}900 trials. Our results reveal that defenses relying on a single observable signal---including input encoding, attention probing, and architecture-level containment---are defeated by adaptive attack (attack success rate rises from 0.73--0.88 to 1.00 on GPT-4o-mini), while runtime-checking defenses resist only by over-blocking benign actions (false positive rate 0.47--1.00). We show this defeat is backbone-dependent: on GPT-4o, input-encoding defenses partially resist (R-ASR drops to 0.13) because the stronger backbone more faithfully obeys delimiter instructions, whereas attention-probing defenses fail on both backbones. We formalize a single-signal fragility conjecture, identify four root-cause families that cluster by defense class, and instantiate four design-principle prototypes, finding that orthogonal multi-signal defenses are necessary but incur a utility cost. We release our framework and benchmark as open-source artifacts.
\end{abstract}
